% GENERATED FILE. Edit canonical lessons, then run npm run generate:books.
% canonical-source-hash: fnv1a64:a3880b44bfc5f186
% canonical-lessons: TA-C321-catanku, TA-C321-marapu, TA-C321-peyarcuttu, TA-C321-valaikappu, TA-C321-putumanai-pukuvila

\chapter{Customs and Rituals}
\label{ch:ta-customs-and-rituals}
\hlchaptermodality{321}

\begin{chapteropening}
\textbf{By the end of this chapter:} \emph{I can talk about family customs and ceremonies.}

Complete the last lesson of chapter 321: I can talk about family customs and ceremonies.
\end{chapteropening}

\section[\texorpdfstring{\ta{சடங்கு} (caṭaṅku)}{சடங்கு (caṭaṅku)}]{\ta{சடங்கு} (caṭaṅku) — a ritual; a rite}

\label{lesson:TA-C321-catanku}



\begin{quote}
Before the new one: say the Tamil for a banner, then the Tamil for to proclaim.
\end{quote}

\subsection*{You'll want to know: \ta{சடங்கு}}
\textbf{\ta{சடங்கு}} — \emph{caṭaṅku} — \textquotedblleft{}a ritual; a rite\textquotedblright{}.

\begin{quote}
\ta{திருமணத்தில்} \ta{பல} \ta{சடங்குகள்} \ta{இருந்தன}.
\end{quote}

\begin{quote}
\emph{tirumaṇattil pala caṭaṅkukaḷ iruntaṉa.}
\end{quote}

\begin{quote}
There were many rituals at the wedding.
\end{quote}

Plural: \ta{சடங்குகள்}, \emph{caṭaṅkukaḷ}.

Plural \ta{சடங்குகள்}.

\begin{cousinweb}
\textbf{\ta{சடங்கு}} (\emph{caṭaṅku}), \textquotedblleft{}a rite, a ceremony\textquotedblright{}, is a Tamil word whose origin is not settled in the sources at hand.
\end{cousinweb}

\begin{grammarlens}[title={What you've built}]
One more word for telling what happened.
\end{grammarlens}

\subsection*{Your turn}
\begin{itemize}
\raggedright
  \item \emph{Say it:} \emph{caṭaṅku}
  \item \emph{Say it:} \emph{caṭaṅku}, once more
  \item \emph{Say it:} \emph{tirumaṇattil pala caṭaṅkukaḷ iruntaṉa.} — There were many rituals at the wedding.
  \item \emph{Recall:} say the Tamil for a banner, then the Tamil for to proclaim, then say \emph{caṭaṅku} again
\end{itemize}

\begin{tcolorbox}[breakable,colback=teal!4,colframe=teal!35!black,title={Before you move on}]
What does \ta{சடங்கு} mean? (\textquotedblleft{}A ritual; a rite\textquotedblright{}.) Say it once more.
\end{tcolorbox}
\section[\texorpdfstring{\ta{மரபு} (marapu)}{மரபு (marapu)}]{\ta{மரபு} (marapu) — a convention; a tradition}

\label{lesson:TA-C321-marapu}



\begin{quote}
Before the new one: say the Tamil for to proclaim, then the Tamil for a ritual.
\end{quote}

\subsection*{You'll want to know: \ta{மரபு}}
\textbf{\ta{மரபு}} — \emph{marapu} — \textquotedblleft{}a convention; a tradition\textquotedblright{}.

\begin{quote}
\ta{இது} \ta{எங்கள்} \ta{குடும்பத்தின்} \ta{பழைய} \ta{மரபு}.
\end{quote}

\begin{quote}
\emph{itu eṅkaḷ kuṭumpattiṉ paḻaiya marapu.}
\end{quote}

\begin{quote}
This is an old convention in our family.
\end{quote}

Also lineage, the line of a family.

\begin{cousinweb}
\textbf{\ta{மரபு}} (\emph{marapu}) is native Dravidian (DEDR 4715): lineage and established usage, what comes down a family line.
\end{cousinweb}

\begin{grammarlens}[title={What you've built}]
One more word for telling what happened.
\end{grammarlens}

\subsection*{Your turn}
\begin{itemize}
\raggedright
  \item \emph{Say it:} \emph{marapu}
  \item \emph{Say it:} \emph{marapu}, once more
  \item \emph{Say it:} \emph{itu eṅkaḷ kuṭumpattiṉ paḻaiya marapu.} — This is an old convention in our family.
  \item \emph{Recall:} say the Tamil for to proclaim, then the Tamil for a ritual, then say \emph{marapu} again
\end{itemize}

\begin{tcolorbox}[breakable,colback=teal!4,colframe=teal!35!black,title={Before you move on}]
What does \ta{மரபு} mean? (\textquotedblleft{}A convention; a tradition\textquotedblright{}.) Say it once more.
\end{tcolorbox}
\section[\texorpdfstring{\ta{பெயர்சூட்டு} (peyarcūṭṭu)}{பெயர்சூட்டு (peyarcūṭṭu)}]{\ta{பெயர்சூட்டு} (peyarcūṭṭu) — to name a child}

\label{lesson:TA-C321-peyarcuttu}



\begin{quote}
Before the new one: say the Tamil for a ritual, then the Tamil for a convention.
\end{quote}

\subsection*{You'll want to know: \ta{பெயர்சூட்டு}}
\textbf{\ta{பெயர்சூட்டு}} — \emph{peyarcūṭṭu} — \textquotedblleft{}to name a child\textquotedblright{}.

\begin{quote}
\ta{பாட்டி} \ta{குழந்தைக்குப்} \ta{பெயர்சூட்டினார்}.
\end{quote}

\begin{quote}
\emph{pāṭṭi kuḻantaikkup peyarcūṭṭiṉār.}
\end{quote}

\begin{quote}
Grandmother named the baby.
\end{quote}

Past (I …): \ta{பெயர்சூட்டினேன்}, \emph{peyarcūṭṭiṉēṉ}.

The child takes -\ta{க்கு}: \ta{குழந்தைக்குப்} \ta{பெயர்சூட்டு}.

\begin{cousinweb}
\textbf{\ta{பெயர்சூட்டு}} (\emph{peyarcūṭṭu}) is a Tamil compound built from \ta{பெயர்}, \textquotedblleft{}name\textquotedblright{} (DEDR 4410), and \ta{சூட்டு}, \textquotedblleft{}to adorn, crown, invest with\textquotedblright{} (DEDR 2721).
\end{cousinweb}

\begin{grammarlens}[title={What you've built}]
One more word for telling what happened.
\end{grammarlens}

\subsection*{Your turn}
\begin{itemize}
\raggedright
  \item \emph{Say it:} \emph{peyarcūṭṭu}
  \item \emph{Say it:} \emph{peyarcūṭṭu}, once more
  \item \emph{Say it:} \emph{pāṭṭi kuḻantaikkup peyarcūṭṭiṉār.} — Grandmother named the baby.
  \item \emph{Recall:} say the Tamil for a ritual, then the Tamil for a convention, then say \emph{peyarcūṭṭu} again
\end{itemize}

\begin{tcolorbox}[breakable,colback=teal!4,colframe=teal!35!black,title={Before you move on}]
What does \ta{பெயர்சூட்டு} mean? (\textquotedblleft{}To name a child\textquotedblright{}.) Say it once more.
\end{tcolorbox}
\section[\texorpdfstring{\ta{வளைகாப்பு} (vaḷaikāppu)}{வளைகாப்பு (vaḷaikāppu)}]{\ta{வளைகாப்பு} (vaḷaikāppu) — a baby shower}

\label{lesson:TA-C321-valaikappu}



\begin{quote}
Before the new one: say the Tamil for a convention, then the Tamil for to name a child.
\end{quote}

\subsection*{You'll want to know: \ta{வளைகாப்பு}}
\textbf{\ta{வளைகாப்பு}} — \emph{vaḷaikāppu} — \textquotedblleft{}a baby shower\textquotedblright{}.

\begin{quote}
\ta{என்} \ta{அக்காவுக்கு} \ta{நேற்று} \ta{வளைகாப்பு} \ta{நடந்தது}.
\end{quote}

\begin{quote}
\emph{eṉ akkāvukku nēṟṟu vaḷaikāppu naṭantatu.}
\end{quote}

\begin{quote}
My elder sister had her baby shower yesterday.
\end{quote}

\ta{வளை}, bangle + \ta{காப்பு}, protection: bangles are put on the mother-to-be.

\begin{cousinweb}
\textbf{\ta{வளைகாப்பு}} (\emph{vaḷaikāppu}) is a Tamil compound built from \ta{வளை}, \textquotedblleft{}a bangle\textquotedblright{}, from the root of \textquotedblleft{}circle\textquotedblright{} (DEDR 5313), and \ta{காப்பு}, \textquotedblleft{}protection\textquotedblright{} (DEDR 1416).
\end{cousinweb}

\begin{grammarlens}[title={What you've built}]
One more word for telling what happened.
\end{grammarlens}

\subsection*{Your turn}
\begin{itemize}
\raggedright
  \item \emph{Say it:} \emph{vaḷaikāppu}
  \item \emph{Say it:} \emph{vaḷaikāppu}, once more
  \item \emph{Say it:} \emph{eṉ akkāvukku nēṟṟu vaḷaikāppu naṭantatu.} — My elder sister had her baby shower yesterday.
  \item \emph{Recall:} say the Tamil for a convention, then the Tamil for to name a child, then say \emph{vaḷaikāppu} again
\end{itemize}

\begin{tcolorbox}[breakable,colback=teal!4,colframe=teal!35!black,title={Before you move on}]
What does \ta{வளைகாப்பு} mean? (\textquotedblleft{}A baby shower\textquotedblright{}.) Say it once more.
\end{tcolorbox}
\section[\texorpdfstring{\ta{புதுமனை} \ta{புகுவிழா} (putumaṉai pukuviḻā)}{புதுமனை புகுவிழா (putumaṉai pukuviḻā)}]{\ta{புதுமனை} \ta{புகுவிழா} (putumaṉai pukuviḻā) — a housewarming}

\label{lesson:TA-C321-putumanai-pukuvila}



\begin{quote}
Before the new one: say the Tamil for to name a child, then the Tamil for a baby shower.
\end{quote}

\subsection*{You'll want to know: \ta{புதுமனை} \ta{புகுவிழா}}
\textbf{\ta{புதுமனை} \ta{புகுவிழா}} — \emph{putumaṉai pukuviḻā} — \textquotedblleft{}a housewarming\textquotedblright{}.

\begin{quote}
\ta{நண்பர்} \ta{என்னைத்} \ta{தன்} \ta{புதுமனை} \ta{புகுவிழாவுக்கு} \ta{அழைத்தார்}.
\end{quote}

\begin{quote}
\emph{naṇpar eṉṉait taṉ putumaṉai pukuviḻāvukku aḻaittār.}
\end{quote}

\begin{quote}
My friend invited me to his housewarming.
\end{quote}

\ta{புது} \ta{மனை}, new house + \ta{புகு}, to enter + \ta{விழா}.

\begin{cousinweb}
\textbf{\ta{புதுமனை} \ta{புகுவிழா}} (\emph{putumaṉai pukuviḻā}) is a Tamil phrase built from \ta{புது}, \textquotedblleft{}new\textquotedblright{} (DEDR 4275), \ta{மனை}, \textquotedblleft{}house\textquotedblright{} (DEDR 4776), \ta{புகு}, \textquotedblleft{}to enter\textquotedblright{} (DEDR 4238), and \ta{விழா}, \textquotedblleft{}festival\textquotedblright{}.
\end{cousinweb}

\begin{grammarlens}[title={What you've built}]
One more word for telling what happened.
\end{grammarlens}

\subsection*{Your turn}
\begin{itemize}
\raggedright
  \item \emph{Say it:} \emph{putumaṉai pukuviḻā}
  \item \emph{Say it:} \emph{putumaṉai pukuviḻā}, once more
  \item \emph{Say it:} \emph{naṇpar eṉṉait taṉ putumaṉai pukuviḻāvukku aḻaittār.} — My friend invited me to his housewarming.
  \item \emph{Recall:} say the Tamil for to name a child, then the Tamil for a baby shower, then say \emph{putumaṉai pukuviḻā} again
\end{itemize}

\begin{tcolorbox}[breakable,colback=teal!4,colframe=teal!35!black,title={Before you move on}]
What does \ta{புதுமனை} \ta{புகுவிழா} mean? (\textquotedblleft{}A housewarming\textquotedblright{}.) Say it once more.
\end{tcolorbox}
